%======================================================================
\section{An \texorpdfstring{\(r^2\)}{r\textasciicircum 2}-weighted apex coefficient (new in v6)}
\label{sec:weighted}
%======================================================================

Constant transverse profiles see only \(\Ccb_b\).
A quadratic transverse jet \(G=r^2g\) needs the next moment of the
apex heat distribution.
v5 recorded that no value of the higher coefficient \(\mathcal D_b\)
had been supplied or derived.
This section derives it with the same Abel subtraction used in
Definition~\ref{def:apex}.

\begin{definition}[Weighted Abel apex trace]
\label{def:weighted}
For \(u>0\) put
\begin{align*}
T^{(2)}_b(u)
 &=\int_{C_\alpha}r^2
 e^{-(\cosh u-1)r^2/(2s)}K_b^\Fr(s;x,x)\,\dd V_\perp,\\
T^{(2)}_{\rm bulk}(u)
 &=\int_{C_\alpha}r^2
 e^{-(\cosh u-1)r^2/(2s)}\frac{\dd V_\perp}{4\pi s}
 =\frac{\alpha s}{(\cosh u-1)^2}.
\end{align*}
When the limit exists, set
\(\ApTr[r^2e^{-sA_b}]=\lim_{u\downarrow0}
(T^{(2)}_b(u)-T^{(2)}_{\rm bulk}(u))\).
\end{definition}

\begin{theorem}[Weighted twisted Friedrichs apex coefficient]
\label{thm:weighted}
Let \(\alpha>0\), \(0\le b<1\), and \(s>0\).
With Bernoulli polynomials \(B_2(b)=b^2-b+\tfrac16\) and
\(B_4(b)=b^4-2b^3+b^2-\tfrac1{30}\),
\begin{equation}
\ApTr\bigl[r^2e^{-sA_b^\Fr}\bigr]=\mathcal D_b\,s,
\qquad
\mathcal D_b=-\frac{11\alpha}{180}
+\frac{B_2(b)}{3\alpha}-\frac{B_4(b)}{6\alpha^3}.
\label{eq:Db}
\end{equation}
Equivalently,
\begin{equation}
\mathcal D_b=
\frac{1+10\alpha^2-11\alpha^4-60\alpha^2b(1-b)-30b^2(1-b)^2}
{180\alpha^3}.
\label{eq:Dbexpanded}
\end{equation}
For a heat superposition \(f\) as in Theorem~\ref{thm:apex},
\begin{equation}
\ApTr\bigl[r^2f(A_b^\Fr)\bigr]=-\mathcal D_b\,f'(0).
\label{eq:Dfunctional}
\end{equation}
\end{theorem}

\begin{proof}
With \(z=r^2/(2s)\), we have \(r^2=2sz\) and
\(r\,\dd r=s\,\dd z\).
The angular integral and the heat representation
\eqref{eq:heatkernel} then give
\[
T^{(2)}_b(u)=s\sum_{n\in\ZZ}\int_0^\infty
z\,e^{-z\cosh u}I_{\abs{n+b}/\alpha}(z)\,\dd z.
\]
Differentiate the Laplace transform from the proof of
Theorem~\ref{thm:apex} with respect to \(u\):
\[
\int_0^\infty z\,e^{-z\cosh u}I_\rho(z)\,\dd z
=-\frac1{\sinh u}\frac{\dd}{\dd u}
\frac{e^{-\rho u}}{\sinh u},
\qquad \rho>-1.
\]
Since all summands are nonnegative, the sum may be interchanged
with the integral, giving
\[
\frac{T^{(2)}_b(u)}{s}=-\frac1{\sinh u}\frac{\dd}{\dd u}
\left[\frac{\cosh((\tfrac12-b)u/\alpha)}
{\sinh u\,\sinh(u/(2\alpha))}\right].
\]
Laurent expansion at \(u=0\) shows that the
\(u^{-4},u^{-3},u^{-2},u^{-1}\) terms of
\(T^{(2)}_b/s\) and \(T^{(2)}_{\rm bulk}/s=\alpha/(\cosh u-1)^2\)
coincide.
The difference of their constant terms is \eqref{eq:Db}.
Expanding the Bernoulli polynomials gives \eqref{eq:Dbexpanded}.

For a heat superposition \(f(\lambda)=\int w(s)e^{-s\lambda}\dd s\),
dominated convergence on the compact proper-time support gives
\(\ApTr[r^2f(A)]=\mathcal D_b\int s\,w(s)\,\dd s
=-\mathcal D_bf'(0)\).
\textsc{[exec]} The Laurent algebra is carried out in
\loc{code/weighted_apex.py}: all four pole terms vanish identically,
and the constant term equals \eqref{eq:Db}
(\loc{logs/v6/weighted_checks.json}, key \texttt{D\_matches\_log}).
\end{proof}

\begin{proposition}[Weighted singular-channel shift]
\label{prop:weightedflip}
For \(0<\nu<1\), replacing one regular channel of order \(\nu\) by
its singular realization changes the weighted apex coefficient by
\[
\Delta\mathcal D=\frac{2\nu(1-\nu^2)}{3}.
\]
For the CR weight-two channel, \(\nu=2-\alpha^{-1}\) and
\[
\mathcal D^\CR_{S=2}=\mathcal D_{2\alpha-1}+\Delta\mathcal D
=\mathcal D_{b=2\alpha},
\]
the polynomial \eqref{eq:Dbexpanded} evaluated at \(b=2\alpha\).
\end{proposition}

\begin{proof}
The weighted analogue of the proof of Theorem~\ref{thm:flip}
replaces \(e^{-\nu u}\) by \(e^{+\nu u}\) in one channel.
The difference is
\(-(\sinh u)^{-1}\partial_u[2\sinh(\nu u)/\sinh u]\).
Its limit as \(u\to0\) is
\(-u^{-1}\partial_u[2\nu(1+(\nu^2-1)u^2/6)]
\to2\nu(1-\nu^2)/3\).
Near zero the singular kernel is \(O(z^{1-\nu})\), so the weighted
diagonal difference is absolutely integrable.
\textsc{[exec]} The identity
\(\mathcal D_{2\alpha}-\mathcal D_{2\alpha-1}=\Delta\mathcal D\) is
checked exactly as a rational function of \(\alpha\)
(\loc{logs/v6/weighted_checks.json}).
\end{proof}

\begin{corollary}[Comparison-sector values]
\label{cor:Dvalues}
For \(1/2<\alpha<1\):
\begin{align*}
\mathcal D_0&=\frac{(1-\alpha^2)(1+11\alpha^2)}{180\alpha^3},\\
\mathcal D_{S=1}&=\frac{(\alpha-1)(\alpha+1)(19\alpha^2-1)}{180\alpha^3},\\
\mathcal D^\Fr_{S=2}&=-\frac{(\alpha-1)(251\alpha^3-829\alpha^2+601\alpha-119)}{180\alpha^3},\\
\mathcal D^\CR_{S=2}&=-\frac{(\alpha-1)(251\alpha^3-109\alpha^2+\alpha+1)}{180\alpha^3}.
\end{align*}
All four vanish as \(\alpha\uparrow1\).
In particular, the weighted CR coefficient has \emph{no} flat
remnant, unlike \(c_2^\CR\to1\) (Remark~\ref{rem:flatCR}),
because \(\Delta\mathcal D\to0\) as \(\nu\to1\).
\end{corollary}

\begin{proposition}[Orbifold checks]
\label{prop:orbifold}
For \(\alpha=1/N\) with \(N\in\{2,\dots,8\}\),
\[
\mathcal D_0(1/N)=\frac1{4N}\sum_{k=1}^{N-1}\csc^4\frac{\pi k}{N}
=\frac{(N^2-1)(N^2+11)}{180N}.
\]
This gives \(1/8,\,8/27,\,9/16,\,24/25,\,329/216,\,16/7,\,105/32\).
For the twisted orbifolds \((N,j)=(3,1),(4,1),(5,2)\), with
\(b=j/N\), the image sums weighted by holonomy phases equal
\(\mathcal D_b\), namely \(-4/27\), \(-1/16\), and \(-16/25\).
\end{proposition}

\begin{proof}
On \(\RR^2/\ZZ_N\) the heat kernel is the image sum.
The \(k\)-th image contributes
\((4\pi s)^{-1}e^{-r^2\sin^2(\pi k/N)/s}\).
Integrating \(r^2\) times it over the wedge of angle \(2\pi/N\) gives
\(s/(4N\sin^4(\pi k/N))\).
The identity element is the bulk term.
\textsc{[exec]} Both sides are evaluated exactly for each \(N\),
and to \(10^{-12}\) for the twisted cases.
See the key \texttt{orbifold\_csc4} in
\loc{logs/v6/weighted_checks.json}, and
\loc{logs/v6/weighted_apex.json}.
\end{proof}

\begin{remark}[Literature comparison]
The untwisted moment \(\mathcal D_0\) has the
\(\csc^4\)-sum structure of the classical conical heat coefficients
\cite{Fursaev1994}.
The identification of \(2\pi\alpha\,\mathcal D_0\) with the
\(C_4\)-type coefficient of that reference is made from memory of its
normalization (\(\dagger\)). It is not used as an input: the
orbifold identity above is self-contained.

Two literature results are analogues of the CR one-sided remnant
of \(c_2^\CR\): the non-uniform spin-dependent flat limits of
\cite{FursaevMiele1997}, and the spin-one horizon contact term of
\cite{Kabat}.
Corollary~\ref{cor:Dvalues} shows that this remnant does
\emph{not} propagate to the \(r^2\) moment.
No numerical spin-2 coefficient is imported from either work.
\end{remark}

\subsection{Application: the second \texorpdfstring{\(\phi\phi\)}{phi-phi} combination (partial)}

By \eqref{eq:Dfunctional}, a quadratic transverse jet contributes
\(-\mathcal D_S\) times the derivative, at zero transverse spectral
parameter, of the longitudinal integrand.
Within Assumption~\ref{ass:reduction}:
\begin{itemize}
\item \textsc{[exec]} \emph{D05a.} After the weighted insertion,
the single-line tadpole has longitudinal integrand proportional to
\(p^2\mathcal D_S/l^4\).
At \(D=2\) this is \(O(R^{-4})\) at large radius, so it has no
logarithmic UV term (Lemma~\ref{lem:UV}).
The D05a contribution to the quadratic-profile pole is therefore
\emph{computed} to be \(0\).
\item \emph{D04.} The bubble couples a graviton line and a scalar
line in different channels.
Its \(r^2\) insertion requires mixed-channel traces of the
transverse position operators \(x_\pm\) between channels \(S\) and
\(S\pm1\).
These are not functions of a single \(A_S\), so they fall outside
Theorem~\ref{thm:weighted} and outside (R2).
The D04 part is \(\OPEN\).
\end{itemize}
The second scalar basis value is therefore \emph{partial}.
It is not reported as equal to its D05a part.
